\documentclass[aspectratio=169]{beamer}
\usepackage[english]{babel}
\usepackage[utf8]{inputenc}
\usepackage[T1]{fontenc}
\usepackage{lmodern}
\usepackage{listings}
\usepackage{xcolor}
\usepackage{graphicx}
\usepackage{textcomp}
\DeclareUnicodeCharacter{2014}{---}
\DeclareUnicodeCharacter{2013}{--}
\DeclareUnicodeCharacter{2212}{-}
\DeclareUnicodeCharacter{2192}{\ensuremath{\rightarrow}}
\DeclareUnicodeCharacter{00B1}{\ensuremath{\pm}}
\DeclareUnicodeCharacter{00D7}{\ensuremath{\times}}
\DeclareUnicodeCharacter{00B7}{\ensuremath{\cdot}}
\DeclareUnicodeCharacter{20AC}{EUR}
\DeclareUnicodeCharacter{2070}{\textsuperscript{0}}
\DeclareUnicodeCharacter{00B9}{\textsuperscript{1}}
\DeclareUnicodeCharacter{00B2}{\textsuperscript{2}}
\DeclareUnicodeCharacter{00B3}{\textsuperscript{3}}
\DeclareUnicodeCharacter{2074}{\textsuperscript{4}}
\DeclareUnicodeCharacter{2075}{\textsuperscript{5}}
\DeclareUnicodeCharacter{2076}{\textsuperscript{6}}
\DeclareUnicodeCharacter{2077}{\textsuperscript{7}}
\DeclareUnicodeCharacter{2078}{\textsuperscript{8}}
\DeclareUnicodeCharacter{2079}{\textsuperscript{9}}
\DeclareUnicodeCharacter{2248}{\ensuremath{\approx}}
\DeclareUnicodeCharacter{2264}{\ensuremath{\le}}

% Colors for code
\definecolor{codebg}{RGB}{245,245,245}
\definecolor{codecomment}{RGB}{0,128,0}
\definecolor{codekeyword}{RGB}{0,0,255}
\definecolor{codestring}{RGB}{163,21,21}
\definecolor{bandcolor}{RGB}{200,50,50}

\lstdefinestyle{pythonstyle}{
    language=Python,
    backgroundcolor=\color{codebg},
    basicstyle=\ttfamily\scriptsize,
    keywordstyle=\color{codekeyword},
    commentstyle=\color{codecomment},
    stringstyle=\color{codestring},
    showstringspaces=false,
    breaklines=true,
    frame=single,
    rulecolor=\color{black},
    escapeinside={(*@}{@*)},
    moredelim=**[l][\color{bandcolor}\bfseries]{\#\ ---\ NEW},
    moredelim=**[l][\color{bandcolor}\bfseries]{\#\ ---\ CHANGED},
    literate=
      {á}{{\'a}}1 {é}{{\'e}}1 {í}{{\'i}}1 {ó}{{\'o}}1 {ú}{{\'u}}1
      {Á}{{\'A}}1 {É}{{\'E}}1 {Í}{{\'I}}1 {Ó}{{\'O}}1 {Ú}{{\'U}}1
      {ñ}{{\~n}}1 {Ñ}{{\~N}}1 {ü}{{\"u}}1 {Ü}{{\"U}}1
      {¿}{{?`}}1 {¡}{{!`}}1
      {—}{{---}}1 {–}{{--}}1 {−}{{-}}1
      {±}{{\ensuremath{\pm}}}1 {×}{{\ensuremath{\times}}}1
      {→}{{\ensuremath{\rightarrow}}}1
      {°}{{\textdegree}}1
      {·}{{\ensuremath{\cdot}}}1
      {€}{{\texteuro}}1
      {≈}{{\ensuremath{\approx}}}1
      {≤}{{\ensuremath{\le}}}1
      {⁰}{{\textsuperscript{0}}}1 {¹}{{\textsuperscript{1}}}1
      {²}{{\textsuperscript{2}}}1 {³}{{\textsuperscript{3}}}1
      {⁴}{{\textsuperscript{4}}}1 {⁵}{{\textsuperscript{5}}}1
      {⁶}{{\textsuperscript{6}}}1 {⁷}{{\textsuperscript{7}}}1
      {⁸}{{\textsuperscript{8}}}1 {⁹}{{\textsuperscript{9}}}1
}

\lstset{style=pythonstyle}

% Title slide macro: \lessontitle{Lesson Number}{Title}
\newcommand{\lessontitle}[2]{
    \title{Lesson #1: #2}
    \author{}
    \institute{Tec de Monterrey}
    \begin{frame}[plain]
        \titlepage
    \end{frame}
}

% Figure frame macro: \figureframe{Title}{Image path}
\newcommand{\figureframe}[2]{
    \begin{frame}{#1}
        \begin{center}
            \includegraphics[height=0.8\textheight,width=\textwidth,keepaspectratio]{#2}
        \end{center}
    \end{frame}
}


% Listings pulled from src/ with \lstinputlisting, so the slide is the file.
\lstset{
    basicstyle=\ttfamily\fontsize{8pt}{9.6pt}\selectfont,
    breaklines=true,
    breakatwhitespace=true,
    columns=fullflexible,
    keepspaces=true,
    xleftmargin=0.3em,
    framesep=2pt,
    aboveskip=0pt,
    belowskip=0pt
}

\newcommand{\resultframe}[3]{%
    \begin{frame}{#1}
        \begin{center}
            \includegraphics[height=0.62\textheight,width=\textwidth,keepaspectratio]{#2}
        \end{center}
        \vspace{0.25em}
        {\small #3\par}
    \end{frame}%
}

\newcommand{\claimframe}[2]{%
    \begin{frame}{#1}
        {\small #2\par}
    \end{frame}%
}

\date{}

\begin{document}

\lessontitle{05}{Mutation}

\section{This lesson}

\begin{frame}{This lesson}
Lesson 04 split crossover into two families that cannot be swapped. Mutation is the remaining operator, still measured on one chromosome.

\vspace{0.6em}
\begin{itemize}
    \item Reach is a rate times a step; both have to be visible.
    \item A permutation chromosome cannot take Gaussian noise.
    \item An optimum with no volume is an artefact a search cannot occupy.
\end{itemize}
\end{frame}

% ================= mutation_reach_01_the_instrument =================
\begin{frame}{A mutation kernel}
For gene \(i\),
\[
B_i\sim\operatorname{Bernoulli}(p),\qquad
\Delta_i\sim\mathcal N(\mu,\sigma^2),\qquad
x_i'=\Pi_{[l,u]}(x_i+B_i\Delta_i).
\]
Thus \(\mathbb E[\sum_i B_i]=Lp\): the rate controls how many genes are
activated, while \(\sigma\) controls displacement. Measure changed positions by
\(H=\sum_i\mathbf 1[x_i\ne x_i']\) and permutation travel by
\(D=\sum_v|\operatorname{pos}_\pi(v)-\operatorname{pos}_{\pi'}(v)|\). A zero-width optimum has
probability zero under an ideal continuous proposal, even if a grid contains it.
\end{frame}

\section{The reach of one mutation, and how to measure it}

\begin{frame}{The reach of one mutation, and how to measure it}
Selection discards and crossover recombines; both work only with values the
population already holds. Mutation is the one operator that can produce a value
nobody had. The whole of this lesson is therefore about one question: HOW FAR
can one mutation go?
\end{frame}

\resultframe{The reach of one mutation, and how to measure it}{../figures/mutation_reach_01_the_instrument.png}{%
    Lesson 01's two numbers reappear exactly — sigma 1.0 arrives \textbf{0 of 2000}, sigma 3.0 \textbf{110 of 2000}. Mean travel is 0.798·sigma until the walls start eating it: at sigma 6.0, \textbf{19.4\%} of proposals are clamped and the ratio drops to 0.690. The trade in one table: sigma 0.5 finds the most uphill moves (\textbf{272}) and arrives \textbf{0} times}

% ================= mutation_reach_02_rate_and_step =================
\section{Mutation has two dials, not one}

\begin{frame}{Mutation has two dials, not one}
Step 1 measured a single gene, so sigma was the only thing there was to set.
A real chromosome has many genes, and the operator then has a second dial: the
per-gene probability p that a gene is touched at all. Gridin's
\texttt{mutation\_random\_deviation(ind, mu, sigma, p)} carries both.
\end{frame}

\begin{frame}[fragile]{(1) GENE\_COUNT}
\lstinputlisting[basicstyle=\ttfamily\fontsize{8pt}{9.6pt}\selectfont,firstline=37,lastline=43]{../src/mutation_reach_02_rate_and_step.py}
\end{frame}

\begin{frame}[fragile]{(2) Individual.fitness}
\lstinputlisting[basicstyle=\ttfamily\fontsize{8pt}{9.6pt}\selectfont,firstline=65,lastline=65]{../src/mutation_reach_02_rate_and_step.py}
\end{frame}

\begin{frame}[fragile]{(3) mutate\_random\_deviation()}
\lstinputlisting[basicstyle=\ttfamily\fontsize{8pt}{9.6pt}\selectfont,firstline=81,lastline=96]{../src/mutation_reach_02_rate_and_step.py}
\end{frame}

\begin{frame}[fragile]{(4) regime\_report()}
\lstinputlisting[basicstyle=\ttfamily\fontsize{6pt}{7.2pt}\selectfont,firstline=99,lastline=135]{../src/mutation_reach_02_rate_and_step.py}
\end{frame}

\resultframe{Mutation has two dials, not one}{../figures/mutation_reach_02_rate_and_step.png}{%
    Genes touched never differs from \texttt{p·n} by more than \textbf{0.015}. Three regimes with the same \texttt{p·sigma = 0.30} travel 1.904 / 2.179 / 2.304 in total but their longest single step differs by a factor of \textbf{7.1} (10.253 vs 1.449). At p = 0.1, \textbf{718 of 2000} mutations change nothing at all, against a predicted (1−p)¹⁰ = 0.3487}

% ================= mutation_reach_03_fitness_driven =================
\section{Mutation that refuses to make things worse}

\begin{frame}{Mutation that refuses to make things worse}
Every mutation so far has been blind: it moves the genes and accepts whatever
the objective says about the result. Most of the time the result is worse -
step 1 already showed that at sigma = 3.0 only 144 of 2000 mutations were
uphill. Gridin's last operator refuses those: it retries, up to MAX\_TRIES
times, and returns the first mutant that beats its parent.
\end{frame}

\begin{frame}[fragile]{(1) MAX\_TRIES}
\lstinputlisting[basicstyle=\ttfamily\fontsize{8pt}{9.6pt}\selectfont,firstline=39,lastline=44]{../src/mutation_reach_03_fitness_driven.py}
\end{frame}

\begin{frame}[fragile]{(2) mutate\_fitness\_driven()}
\lstinputlisting[basicstyle=\ttfamily\fontsize{7pt}{8.5pt}\selectfont,firstline=95,lastline=119]{../src/mutation_reach_03_fitness_driven.py}
\end{frame}

\begin{frame}[fragile]{(3) compare()}
\lstinputlisting[basicstyle=\ttfamily\fontsize{6pt}{7.2pt}\selectfont,firstline=122,lastline=165]{../src/mutation_reach_03_fitness_driven.py}
\end{frame}

\resultframe{Mutation that refuses to make things worse}{../figures/mutation_reach_03_fitness_driven.png}{%
    The filter turns a mean fitness change of \textbf{−1.0337 into +0.0548} and doubles arrivals (102 → 273) — but spends \textbf{5,603 fitness calls against 2,000}, so per thousand calls it arrives \textbf{48.7} times against blind mutation's \textbf{51.0}. Acceptance 20.2\% is exactly 1−(1−q)³ for the blind uphill rate q}

% ================= mutation_reach_04_the_dead_staircase =================
\section{The run Lesson 02 left for dead}

\begin{frame}{The run Lesson 02 left for dead}
An earlier version of Lesson 02 step 5 ended with a run that stopped six
generations in, 0.2000 short of the staircase's optimum, with almost no gene
spread left, and pointed here: mutation "is the one operator that could put it
back". This step is where that promise came due - and it did not come due the way
it was written. What the investigation found instead sent a correction back to
Lesson 02, which no longer makes that claim.
\end{frame}

\begin{frame}[fragile]{(1) run()}
\lstinputlisting[basicstyle=\ttfamily\fontsize{6pt}{7.2pt}\selectfont,firstline=115,lastline=201]{../src/mutation_reach_04_the_dead_staircase.py}
\end{frame}

\begin{frame}[fragile]{(2) staircase()}
\lstinputlisting[basicstyle=\ttfamily\fontsize{8pt}{9.6pt}\selectfont,firstline=93,lastline=107]{../src/mutation_reach_04_the_dead_staircase.py}
\end{frame}

\begin{frame}[fragile]{(3) step\_widths()}
\lstinputlisting[basicstyle=\ttfamily\fontsize{7pt}{8.5pt}\selectfont,firstline=204,lastline=223]{../src/mutation_reach_04_the_dead_staircase.py}
\end{frame}

\begin{frame}[fragile]{(4) PEAK\_HEIGHT\_REPAIRED}
\lstinputlisting[basicstyle=\ttfamily\fontsize{8pt}{9.6pt}\selectfont,firstline=226,lastline=233]{../src/mutation_reach_04_the_dead_staircase.py}
\end{frame}

\begin{frame}[fragile]{(5) sweep()}
\lstinputlisting[basicstyle=\ttfamily\fontsize{6pt}{7.2pt}\selectfont,firstline=236,lastline=264]{../src/mutation_reach_04_the_dead_staircase.py}
\end{frame}

\resultframe{The run Lesson 02 left for dead}{../figures/mutation_reach_04_the_dead_staircase.png}{%
    Lesson 02's run reproduced (6 generations, best \textbf{+1.8000}, spread 0.192, \textbf{0.2000} short). Then eight mutation regimes × 100 runs: the top step is reached \textbf{0 times in 800 runs}, while mean gene spread goes from 0.462 to 3.194. \texttt{step\_widths()} says why — the top step is \textbf{0.000 wide}. Repaired (peak height 10.0 → 10.5, top step 0.500 wide), Lesson 02's own regime reaches it \textbf{96 of 100} times}

% ================= permutation_mutation_01_illegal_noise =================
\section{permutation mutation 01 illegal noise}

\begin{frame}{permutation mutation 01 illegal noise}
Lesson 05 - Step 1 of the second example: when noise is not a legal move
=========================================================================
Every operator so far has added a number to a gene. That works because a gene
was a real value and any real value in [-10, +10] was a legal gene.
\end{frame}

\resultframe{permutation mutation 01 illegal noise}{../figures/permutation_mutation_01_illegal_noise.png}{%
    Random deviation on a permutation produces \textbf{1935 invalid mutants of 2000}, and \textbf{0} of the 65 survivors had any gene touched — the operator's only legal output is the untouched one, at a measured rate matching (1−p)¹⁰ to within \textbf{0.0025}. Bit flip changes exactly 1 position every time; exchange exactly 2}

% ================= permutation_mutation_02_rearrangements =================
\section{permutation mutation 02 rearrangements}

\begin{frame}{permutation mutation 02 rearrangements}
Lesson 05 - Step 2 of the second example: three bigger rearrangements
======================================================================
NEW IN THIS STEP: inversion, shift and shuffle, and the comparison that puts
all four permutation operators on one scale.
\end{frame}

\begin{frame}[fragile]{(1) mutation\_inversion()}
\lstinputlisting[basicstyle=\ttfamily\fontsize{7pt}{8.5pt}\selectfont,firstline=79,lastline=97]{../src/permutation_mutation_02_rearrangements.py}
\end{frame}

\begin{frame}[fragile]{(2) mutation\_shift()}
\lstinputlisting[basicstyle=\ttfamily\fontsize{7pt}{8.5pt}\selectfont,firstline=100,lastline=119]{../src/permutation_mutation_02_rearrangements.py}
\end{frame}

\begin{frame}[fragile]{(3) mutation\_shuffle()}
\lstinputlisting[basicstyle=\ttfamily\fontsize{7pt}{8.5pt}\selectfont,firstline=122,lastline=140]{../src/permutation_mutation_02_rearrangements.py}
\end{frame}

\begin{frame}[fragile]{(4) compare\_operators()}
\lstinputlisting[basicstyle=\ttfamily\fontsize{6pt}{7.2pt}\selectfont,firstline=167,lastline=199]{../src/permutation_mutation_02_rearrangements.py}
\end{frame}

\resultframe{permutation mutation 02 rearrangements}{../figures/permutation_mutation_02_rearrangements.png}{%
    The two measurements of reach rank the operators differently: by positions disturbed \textbf{shift (4.663) > inversion (4.230) > shuffle (3.708) > exchange (2.000)}, by distance travelled \textbf{inversion (1.308) > shuffle (0.876) > exchange (0.733) = shift (0.733)}. Exchange and shift give identical total travel on \textbf{90 of 90} position pairs. Shuffle returns the individual unchanged \textbf{271 of 2000} times (13.6\%)}

\section{Next}

\begin{frame}{Next}
Lesson 02's staircase once had a top step of width zero. This lesson keeps the broken geometry as the specimen it diagnoses.

\vspace{0.8em}
\textbf{Lesson 06 --- Measuring Effectiveness}

A genetic algorithm is a random variable. One generation, and one run, prove nothing. Lesson 03 postponed that evidence here.
\end{frame}
\end{document}
